% TikZ diagram for the given BLIF
% - Oriented horizontally: inputs on the left, outputs on the right
% - Each .names block is shown as a small LUT/AND/INV-style logic box
% - Wires follow the BLIF net dependencies

\begin{tikzpicture}[
  x=1.35cm, y=0.9cm,
  >=Latex,
  wire/.style={-Latex, line width=0.5pt},
  nodebox/.style={draw, rounded corners=2pt, align=center, inner sep=3pt, minimum height=8mm},
  port/.style={draw, circle, inner sep=1.2pt},
  lab/.style={font=\small}
]

% Inputs (left)
\node[lab] (x0) at (-1,  3) {$x_0$};
\node[lab] (x1) at (-1,  1.5) {$x_1$};
\node[lab] (x2) at (-1,  0) {$x_2$};
\node[lab] (x3) at (-1, -1.5) {$x_3$};
\node[lab] (x4) at (-1, -3) {$x_4$};

% Internal logic nodes (placed left-to-right following dataflow)

% .names x0 x1 new_n8   (truth table: 00 1)  => new_n8 = ~x0 & ~x1
\node[nodebox] (n8) at (1, 3) {\scriptsize $i_0 = !x_0  !x_1$};
\node[nodebox] (n8b) at (1, 1.5) {\scriptsize $i_1 = !x_0  !x_1$};
% .names x2 new_n8 new_n9  (11 1) => new_n9 = x2 & new_n8
\node[nodebox] (n9) at (2.5, 1) {\scriptsize $i_2 = x_2  i_0$};

% .names x3 new_n9 po0  (01 1) => po0 = ~x3 & new_n9
\node[nodebox] (po0box) at (4, 0) {\scriptsize $po_0 = !x_3  i_2$};

% .names x2 new_n8 new_n11 (01 1) => new_n11 = ~x2 & new_n8
\node[nodebox] (n11) at (2.5, -1) {\scriptsize $ i_3 = !x_2  i_1$};

% .names x3 new_n11 new_n12 (11 1) => new_n12 = x3 & new_n11
\node[nodebox] (n12) at (4, -2) {\scriptsize $i_4 = x_2  i_3$};

% .names x4 new_n12 po1 (01 1) => po1 = ~x4 & new_n12
\node[nodebox] (po1box) at (5.5, -3) {\scriptsize $po_1 = !x_4  i_4$};

% Wires from inputs to logic
\draw[wire] (x0.east) -- (n8.west);
\draw[wire] (x1.east) -- (n8.west);

\draw[wire] (x0.east) -- (n8b.west);
\draw[wire] (x1.east) -- (n8b.west);

\draw[wire] (x2.east) -- (n9.west);
\draw[wire] (n8.east) -- (n9.west);

\draw[wire, bend left=25] (x3.east) to (po0box.west);
\draw[wire] (n9.east)  -- (po0box.west);

\draw[wire] (x2.east) -- (n11.west);
\draw[wire] (n8b.east)  -- (n11.west);

\draw[wire] (x3.east) -- (n12.west);
\draw[wire] (n11.east) -- (n12.west);

\draw[wire] (x4.east) -- (po1box.west);
\draw[wire] (n12.east) -- (po1box.west);

% Output labels on the far right (ports)
\node[lab] ($po_0$) at (10, 0.75) {};
\node[lab] ($po_1$) at (10, -0.75) {};

% Draw explicit output wires to the right (so outputs are clearly on the right edge)
\draw[wire] (po0box.east) -- ++(1.2,0) node[lab, right] {po0};
\draw[wire] (po1box.east) -- ++(1.2,0) node[lab, right] {po1};

\end{tikzpicture}
